\section*{Bab \ref{ch:g11:igneous-rocks} --- Batuan Beku}

\begin{solution}{exo:g11:igneous-rocks:1}
\omterm{def:g11:igneous-rocks:volcanic-plutonic}{Batuan vulkanik} terbentuk dari magma yang membeku di permukaan, sebagai lava;
\omterm{def:g11:igneous-rocks:volcanic-plutonic}{batuan plutonik} mengkristal di kedalaman, di dalam batuan lain. Pasangannya:
basalt dan \omterm{def:g11:ocean-ridges:gabbro}{gabro} (\omterm{def:g11:igneous-rocks:felsic-mafic}{mafik}), \omterm{def:g11:igneous-rocks:rock-names}{andesit} dan \omterm{def:g11:igneous-rocks:rock-names}{diorit} (\omterm{def:g11:igneous-rocks:felsic-mafic}{intermediet}), \omterm{def:g11:igneous-rocks:rock-names}{riolit} dan granit
(\omterm{def:g11:igneous-rocks:felsic-mafic}{felsik}).
\end{solution}

\begin{solution}{exo:g11:igneous-rocks:2}
(1) Gelasan. (2) Berbutir kasar. (3) \omterm{def:g11:igneous-rocks:porphyritic}{Porfiritik} (\omterm{def:g11:igneous-rocks:porphyritic}{fenokris} olivin dalam massa
dasar yang halus). (4) Berbutir halus.
\end{solution}

\begin{solution}{exo:g11:igneous-rocks:3}
\qty{48}{\percent}: \omterm{def:g11:igneous-rocks:felsic-mafic}{mafik}, \omterm{def:g11:ocean-ridges:gabbro}{gabro}. \qty{60}{\percent}: \omterm{def:g11:igneous-rocks:felsic-mafic}{intermediet}, \omterm{def:g11:igneous-rocks:rock-names}{diorit}.
\qty{72}{\percent}: \omterm{def:g11:igneous-rocks:felsic-mafic}{felsik}, granit.
\end{solution}

\begin{solution}{exo:g11:igneous-rocks:4}
\omterm{def:g11:igneous-rocks:intrusion}{Retas} adalah lembaran batuan beku yang memotong lapisan batuan yang
diterobosnya; \omterm{def:g11:igneous-rocks:intrusion}{sill} adalah lembaran yang sejajar dengan lapisan-lapisan, di
antara lapisan-lapisan itu; \omterm{def:g11:igneous-rocks:intrusion}{pluton} adalah \omterm{def:g11:igneous-rocks:intrusion}{intrusi} besar dengan bentuk apa pun
yang terbentuk di kedalaman. \omterm{def:g11:igneous-rocks:intrusion}{Retas} tegak yang naik dari \omterm{def:g11:igneous-rocks:intrusion}{pluton} itulah yang
memasok gunung api.
\end{solution}

\begin{solution}{exo:g11:igneous-rocks:5}
Olivin terbentuk pertama; feldspar kalium, kuarsa, dan kadang-kadang muskovit
terakhir. Plagioklas mula-mula kaya kalsium dan menjadi makin lama makin kaya
natrium ketika magma mendingin.
\end{solution}

\begin{solution}{exo:g11:igneous-rocks:6}
(1) A: \qty{74}{\percent} \ce{SiO2}, \omterm{def:g11:igneous-rocks:felsic-mafic}{felsik}; B: \qty{47}{\percent}, \omterm{def:g11:igneous-rocks:felsic-mafic}{mafik}.
(2) A miskin \ce{FeO}, \ce{MgO}, dan \ce{CaO} serta kaya \ce{Na2O} dan
\ce{K2O}, seperti magma \omterm{def:g11:igneous-rocks:felsic-mafic}{felsik}; B kaya \ce{FeO}, \ce{MgO}, dan \ce{CaO}, serta
miskin \ce{Na2O}, seperti magma \omterm{def:g11:igneous-rocks:felsic-mafic}{mafik}. (3) A adalah \omterm{def:g11:igneous-rocks:rock-names}{riolit}; B \omterm{def:g11:ocean-ridges:gabbro}{gabro}.
\end{solution}

\begin{solution}{exo:g11:igneous-rocks:7}
Waktunya bertambah sebanding dengan kuadrat tebalnya. \qty{1}{m}: $10/4 = 2.5$
hari. \qty{20}{m}: $10 \times 10^2 = 1000$ hari, sekitar 2.7 tahun.
\qty{200}{m}: $10 \times 100^2 = \num{100000}$ hari, sekitar 270 tahun.
\end{solution}

\begin{solution}{exo:g11:igneous-rocks:8}
(1) \unit{m^2} dibagi \unit{m^2/s} menghasilkan detik. (2) $400/10^{-6} =
\qty{4e8}{s}$, yaitu $4 \times 10^8 / 3.16 \times 10^7 \approx 12.7$
tahun. (3) $0.21 \times 12.7 \approx 2.7$ tahun.
\end{solution}

\begin{solution}{exo:g11:igneous-rocks:9}
\omterm{def:g11:igneous-rocks:texture}{Teksturnya} \omterm{def:g11:igneous-rocks:porphyritic}{porfiritik}: \omterm{def:g11:igneous-rocks:porphyritic}{fenokris} plagioklas dalam massa dasar yang berbutir
halus. Mula-mula magma mendingin perlahan di kedalaman, dan kristal-kristal
plagioklas tumbuh besar selagi sebagian besar magma tetap cair; lalu magma
naik dan mendingin dengan cepat, dan sisa zat cairnya mengkristal menjadi
massa dasar yang halus.
\end{solution}

\begin{solution}{exo:g11:igneous-rocks:10}
(1) Kiri: 2 Mg, 2 Si, 6 O; kanan: 2 Mg, 2 Si, 6 O. (2) \qty{1}{kg}
olivin adalah $1000/140.7 \approx 7.11$~mol, yang mengambil mol silika dalam
jumlah yang sama: $7.11 \times 60.1 \approx \qty{427}{g}$. (3) Bila tidak ada
lagi silika dalam zat cair, reaksi itu tidak dapat berlanjut: olivinnya tetap.
\end{solution}

\begin{solution}{exo:g11:igneous-rocks:11}
(1) $2.5 \times 20 \times 1100 = \qty{55000}{MJ}$, \qty{5.5e10}{J}.
(2) $900 \times 20 = \qty{18000}{MJ}$, \qty{1.8e10}{J}. (3) Kalor
kristalisasi kira-kira sepertiga kalor yang hilang karena pendinginan, atau
seperempat dari jumlahnya: cukup untuk sangat memperlambat pembekuan, karena
seluruhnya dilepaskan di atas solidus.
\end{solution}

\begin{solution}{exo:g11:igneous-rocks:12}
(1) Kristal terkecil berada pada bidang kontak, yang membeku dalam satu hari;
yang terbesar di bagian tengah, yang tetap cair selama sekitar 950 hari.
(2) Tidak: \qty{5}{m} pertama memerlukan sekitar 545 hari, \qty{5}{m} terakhir
sekitar 400 hari; mula-mula batas pembekuan itu melambat, lalu bergerak lebih
cepat di dekat bagian tengah, tempat kalor yang keluar melalui kedua sisi
bertemu. (3) Bahkan di bagian tengahnya, \omterm{def:g11:igneous-rocks:intrusion}{sill} itu membeku dalam waktu kurang
dari tiga tahun; granit dengan kristal besar memerlukan waktu yang jauh lebih
lama.
\end{solution}

\begin{solution}{exo:g11:igneous-rocks:13}
(1) \omterm{def:g11:igneous-rocks:intrusion}{Pluton} itu 100 kali lebih tebal: $2.6 \times 100^2 = \num{26000}$
tahun. (2) Misalnya: \omterm{def:g11:igneous-rocks:intrusion}{pluton} yang dalam berada di dalam batuan yang sudah panas,
sehingga kehilangan kalornya lebih lambat daripada ke dalam batuan bersuhu
\qty{50}{\degreeCelsius}; \omterm{def:g11:igneous-rocks:intrusion}{pluton} itu dapat dipasok oleh magma baru berulang
kali; \omterm{def:g11:igneous-rocks:intrusion}{pluton} bukan lembaran yang tak terhingga luasnya; sifat-sifat batuan
berubah dengan suhu.
\end{solution}

\begin{solution}{exo:g11:igneous-rocks:14}
Magma \omterm{def:g11:igneous-rocks:felsic-mafic}{mafik} panas dan miskin silika: magma itu mengkristalkan mineral-mineral
pertama dalam deret --- olivin, piroksen, plagioklas yang kaya kalsium --- dan
kehabisan silika sebelum kuarsa dapat terbentuk. Magma \omterm{def:g11:igneous-rocks:felsic-mafic}{felsik} lebih dingin dan
kaya silika: olivin tidak stabil di dalamnya, dan magma itu berlanjut sampai
akhir deret, membentuk plagioklas yang kaya natrium, feldspar kalium, dan
kuarsa.
\end{solution}

\begin{solution}{exo:g11:igneous-rocks:15}
(1) $b = \sqrt{t\kappa/C} = \sqrt{86400 \times 10^{-6}/0.21} \approx
\qty{0.64}{m}$. (2) Hukum itu memberikan $0.21 \times 25/10^{-6} \approx
\qty{5.3e6}{s}$, sekitar dua bulan: orde besarnya sama dengan satu bulan
(aliran lava juga kehilangan kalor ke udara di bagian atasnya). (3) Gelasan
untuk bahan yang dilontarkan ke udara atau ke air; berbutir halus untuk aliran
lava dan lembaran tipis; lebih kasar, tetapi masih kecil, untuk \omterm{def:g11:igneous-rocks:intrusion}{intrusi} dangkal
yang tebal.
\end{solution}

\begin{solution}{pb:g11:igneous-rocks:1}
\textbf{1.} Sebuah \omterm{def:g11:igneous-rocks:intrusion}{sill}: lembaran itu terletak sejajar dengan lapisan-lapisan,
di antara lapisan-lapisan itu.

\textbf{2.} Pada bidang kontak, sangat berbutir halus, hampir gelasan; di
bagian tengah, berbutir halus sampai sedang, dengan mineral-mineral yang baru
terlihat dengan kaca pembesar.

\textbf{3.} Bidang kontak kehilangan kalornya paling cepat, ke dalam batu pasir
yang dingin: kristal-kristalnya tidak sempat tumbuh.

\textbf{4.} Tepi dingin. Tepi itu membuktikan bahwa magma itu panas ketika
menempati tempatnya, dan batu pasirnya dingin.

\textbf{5.} Sebuah \omterm{def:g11:igneous-rocks:intrusion}{retas}, pemasok \omterm{def:g11:igneous-rocks:intrusion}{sill} itu.

\textbf{6.} \omterm{def:g11:igneous-rocks:felsic-mafic}{Mafik}: \qty{49}{\percent} \ce{SiO2}, kaya akan besi, magnesium, dan
kalsium.

\textbf{7.} Olivin, piroksen, dan plagioklas yang kaya kalsium.

\textbf{8.} Basalt; \omterm{def:g11:igneous-rocks:volcanic-plutonic}{batuan plutonik} yang komposisinya sama adalah \omterm{def:g11:ocean-ridges:gabbro}{gabro}.

\textbf{9.} \omterm{def:g11:igneous-rocks:texture}{Tekstur} bergantung pada laju pendinginan, yang berbeda dari bidang
kontak sampai ke bagian tengah, bukan pada komposisinya.

\textbf{10.} Olivin lebih dahulu; plagioklas mulai kurang lebih pada titik yang
sama dengan piroksen.

\textbf{11.} $b^2/\kappa = 400/10^{-6} = \qty{4e8}{s}$.

\textbf{12.} $0.206 \times 4 \times 10^8 \approx \qty{8.2e7}{s}$, yaitu
$8.2 \times 10^7 / 3.16 \times 10^7 \approx 2.6$ tahun.

\textbf{13.} $0.206 \times 4/10^{-6} \approx \qty{8.2e5}{s}$, sekitar 9.5
hari.

\textbf{14.} $0.07 \times 4 \times 10^8 = \qty{2.8e7}{s}$, sekitar 0.9 tahun.
Magma yang sedang mengkristal melepaskan kalor, yang juga harus dibuang
sebelum magma itu menjadi padat.

\textbf{15.} Sekitar 950 hari dibandingkan dengan 0.6 hari: sekitar 1600 kali
lebih cepat.

\textbf{16.} Sekitar 950 hari, kira-kira 2.6 tahun: jauh lebih lama daripada
aliran lava, di dalam rentang bulan sampai tahun \omterm{def:g11:igneous-rocks:intrusion}{intrusi} dangkal.

\textbf{17.} Waktu seribu kali lebih lama memerlukan tebal $\sqrt{1000} \approx
32$ kali lebih besar: sekitar \qty{630}{m}.

\textbf{18.} Di dekat permukaan, tubuh magma akan meletus, atau mendingin
terlalu cepat; tubuh besar yang mendingin dalam waktu yang sangat lama pasti
berada di kedalaman, di dalam batuan yang sudah panas.

\textbf{19.} Batuan itu telah terkikis habis oleh erosi.

\textbf{20.} \omterm{def:g11:igneous-rocks:intrusion}{Sill} itu memerlukan sekitar \textbf{2.6 tahun} untuk membeku.
\end{solution}
